% Abstract — see issue #27.
% TODO(#27): the results sentence is a placeholder; fill it with the real
% figures once results.tex is written, and check every claim against the body.
\begin{abstract}
Alzheimer's disease (AD) is a progressive neurodegenerative disorder for which early diagnosis is essential, since treatment started in the early stage can slow its progression. Convolutional neural networks (CNNs) applied to structural magnetic resonance imaging (MRI) classify cognitively normal subjects, patients with mild cognitive impairment and patients with AD with high reported accuracy, but their clinical adoption is limited because they behave as black boxes and because many studies split their data by image instead of by subject, which inflates the reported performance. This work proposes an interpretable CNN for AD classification from structural MRI, evaluated with a subject-level split to avoid data leakage. We first review recent classification studies in a comparative matrix and then train and evaluate the network on the three diagnostic classes. Explanation techniques are used to show which brain regions support each prediction and to check whether they agree with the regions known to be affected by the disease. The results report both the classification performance and the interpretability analysis [RESULTS TO BE ADDED]. We conclude that reporting performance under a leakage-free protocol, together with an explanation of each decision, is a requirement for deep learning models to be considered for clinical use.
\end{abstract}
